\section{System Architecture and Methods}
\label{sec:methods}

The method combines the prioritized assignment and route-window procedures of
the centralized baseline \cite{guo2024} with a hierarchical floor decomposition.
The central component assigns jobs, splits routes at lifts, arbitrates lift
use, and merges the resulting schedules. A floor-specific planner constructs
conflict-free timings for its local route legs.

\subsection{Architecture classification}

The architecture is hierarchical and decomposed, with floor-local planning
components; it is not fully decentralized. Job assignment, lift booking,
route-leg dispatch, final-plan merging, and verification remain centrally
orchestrated. Floor planners maintain local node and track reservation sets
$W_n$ and $W'_e$, and resolve conflicts on their own route legs. The central
lift calendar coordinates the cross-floor resource.

Figure~\ref{fig:architecture-flow} summarizes the information flow. The
central-to-floor exchanges are internal planner requests and reservation
responses in the evaluated implementation. The AGVs do not negotiate jobs or
lift slots directly.

\begin{figure}[t]
\centering
\fbox{\parbox{0.94\linewidth}{\centering
Jobs and fleet state $\longrightarrow$ Central planner\\
\emph{Algorithm 1 assignment; route split; lift booking; dispatch}\\
$\longleftrightarrow$ Floor planners $f\in F$\\
\emph{Local route-leg planning with $W_n$ and $W'_e$}\\
$\longrightarrow$ Central merge and independent verification\\
$\longrightarrow$ Complete schedules for AGV execution}}
\caption{Central coordination and floor-local route planning.}
\label{fig:architecture-flow}
\end{figure}

\subsection{Algorithm 1: prioritized job assignment}

Jobs are processed in increasing desired pickup time, with job identifier as
the tie-breaker. For each candidate AGV $i$, the planner evaluates
$F_{ji}=ct_{ji}+\lambda(et_{ji}+dt_{ji})$ using predicted completion time,
pickup earliness, and delivery tardiness. It assigns the job to the AGV with
minimum fitness; ties use fixed AGV priority. The selected AGV's route is
extended through shortest paths to pickup and delivery, with loading and
unloading dwell times $\eta$ and $\eta'$, and is ultimately extended back to
parking $\pi$.

\begin{algorithm}[t]
\caption{Prioritized job assignment}
\label{alg:assignment}
\begin{algorithmic}[1]
\REQUIRE Jobs $J$, AGVs $V$, fitness weight $\lambda$, service times $\eta,\eta'$
\ENSURE Assignment $x_{ji}$ and complete routes $R_i$
\FOR{each $j\in J$ in order $(t_j,j)$}
  \FOR{each $i\in V$}
    \STATE Compute $ct_{ji},et_{ji},dt_{ji}$ from AGV state $(s_i,vt_i)$.
    \STATE $F_{ji}\gets ct_{ji}+\lambda(et_{ji}+dt_{ji})$.
  \ENDFOR
  \STATE Choose $i^*\gets\arg\min_i(F_{ji},\rho_i)$; set $x_{ji^*}\gets1$.
  \STATE Append shortest paths $s_{i^*}\to p_j\to p'_j$ with dwell $\eta,\eta'$.
  \STATE Update AGV state to $(p'_j,ct_{ji^*})$.
\ENDFOR
\FOR{each $i\in V$}
  \STATE Append a shortest path from the final assigned station to $\pi$.
\ENDFOR
\end{algorithmic}
\end{algorithm}

\subsection{Algorithm 2: prioritized conflict-free routing}

Algorithm~\ref{alg:routing} processes AGVs in ascending fixed priority,
$V_1\rightarrow V_{|V|}$. For each route, it constructs node occupancy
windows $[lt_m,rt_m]$ and edge traversal windows. If a window conflicts with a
frozen reservation, Algorithm~\ref{alg:sliding} delays the affected route
segment. The resulting node and track windows are then added to the frozen
sets $W_n$ and $W'_e$. In the paper-style configuration, only parking $\pi$
is an exempt waiting node; the lobby configurations also exempt lift-lobby
nodes.

\begin{algorithm}[t]
\caption{Prioritized conflict-free routing}
\label{alg:routing}
\begin{algorithmic}[1]
\REQUIRE Routes $R_i$, dwell times, priority order $V_1,\ldots,V_{|V|}$
\ENSURE Timed plans and frozen reservations $W_n,W'_e$
\STATE Initialize empty node and edge reservation sets $W_n,W'_e$.
\FOR{each AGV $i$ in priority order}
  \STATE Construct node and edge windows along $R_i$.
  \FOR{each candidate window conflicting with $W_n$ or $W'_e$}
    \STATE Apply Algorithm~\ref{alg:sliding} to delay the route segment.
  \ENDFOR
  \STATE Freeze ordinary node windows in $W_n$ and track windows in $W'_e$.
\ENDFOR
\end{algorithmic}
\end{algorithm}

\subsection{Algorithm 3: sliding-window conflict elimination}

For a candidate interval $[lt,rt]$ and fixed windows $W$, the interval
function $f(lt,rt,W)$ returns the smallest right shift $st$ that clears the
candidate with safety separation $\delta$, together with the available
shift bound $mt$ before the next blocking window. For a conflict at route
visit $p$, the algorithm sweeps backward to the nearest preceding exempt
node $q$, accumulating the route-segment shift $st^*$ and maintaining the
bound $mt^*$. It repeats the sweep up to $k_{\max}$ passes. On exhaustion,
the fallback delays departure at the exempt node until the blocking traffic
is free. The shifted windows from $q+1$ through $p$ move by $st^*$, and the
window at $q$ is extended by that delay.

\begin{algorithm}[t]
\caption{Sliding-window conflict elimination}
\label{alg:sliding}
\begin{algorithmic}[1]
\REQUIRE Conflict visit $p$, reservations $W$, route $R$, limit $k_{\max}$
\ENSURE Conflict-free shifted windows and delay $st^*$
\STATE Find the nearest preceding exempt visit $q$.
\FOR{$k=1$ to $k_{\max}$}
  \STATE Set $st^*\gets0$ and $mt^*\gets\infty$.
  \FOR{visits $m=p,p-1,\ldots,q+1$}
    \STATE Evaluate $f(lt_m+st^*,rt_m+st^*,W)$; update $mt^*$ and $st^*$.
    \IF{$st^*>mt^*$} \STATE Stop this backward sweep. \ENDIF
  \ENDFOR
  \IF{the sweep completes within $mt^*$} \STATE \textbf{break} \ENDIF
\ENDFOR
\IF{no feasible sweep completed}
  \STATE Apply the fallback delay at $q$ until blocking reservations clear.
\ENDIF
\STATE Shift windows $q+1,\ldots,p$ and extend the window at $q$ by $st^*$.
\end{algorithmic}
\end{algorithm}

\subsection{Hybrid central-plus-floor-server extension}

The hybrid planner cuts each complete route at lift edges into floor-local
legs. Each floor planner schedules a leg using local $W_n$ and $W'_e$; when
needed, the AGV waits at its entry lobby or at $\pi$. The central planner
reserves the lift track at the earliest available time after the AGV is
ready, respecting the safety separation. It then releases the next floor leg
at the arrival lobby and merges floor and lift windows into the complete
plan before verification.

Legs are dispatched using either priority ordering or first-come-first-served
(FCFS). FCFS orders legs by release time, AGV priority, then leg index. The
priority order processes the legs of higher-priority AGVs first and is used
as a decomposition check. In the tested seeds, priority-order hybrid plans
match the centralized-with-lobbies plans. For the principal performance
comparison, configuration C uses FCFS order.

\begin{algorithm}[t]
\caption{Hybrid floor-leg planning and lift booking}
\label{alg:hybrid}
\begin{algorithmic}[1]
\REQUIRE Assigned routes, floor planners, central lift calendar
\ENSURE Merged complete AGV plans
\STATE Split each route into maximal floor-local legs at every lift.
\STATE Insert each initial leg into the dispatch queue using the chosen order.
\WHILE{the queue is nonempty}
  \STATE Remove the next leg and schedule it with its floor planner.
  \IF{the route continues on another floor}
    \STATE Reserve the intervening lift trip centrally.
    \STATE Release the next leg at the arrival lobby after lift travel.
  \ENDIF
\ENDWHILE
\STATE Merge floor-leg and lift windows; independently verify each plan.
\end{algorithmic}
\end{algorithm}
